% =========================================================
% capitulos/01_introduccion.tex
% SOLO contenido. Nada de \documentclass ni \begin{document}.
	% Este archivo se llama desde main.tex con \include{...}
	% =========================================================
	
\chapter{Introducción}
\label{chap:introduccion}
	
\section{Contexto y motivación}
	
En el marco del Aprendizaje Máquina (Machine Learning), el objetivo fundamental es el desarrollo de algoritmos capaces de identificar patrones complejos a partir de datos para la toma de decisiones automatizada en diversos campos de aplicación \cite{bharadiya2023tutorial}. En problemas de aprendizaje supervisado, un conjunto de datos puede representarse como una matriz en la que cada fila corresponde a una observación y cada columna a una característica descriptiva. Cuando las observaciones cuentan con una etiqueta de clase, esta representación se convierte en la base sobre la cual los modelos buscan identificar fronteras de separación \cite{hastie2009elements}.
	
Sin embargo, cuando dicha matriz posee un número elevado de características, su análisis y representación gráfica se ven severamente comprometidos debido al fenómeno conocido como la maldición de la dimensionalidad \cite{venkat2018curse}. Para mitigar este problema, las técnicas de reducción de dimensionalidad buscan transformar el espacio original en una representación más compacta, procurando conservar la mayor cantidad de información útil \cite{nguyen2019ten, sorzano2014survey}. 

De manera general, estas técnicas se agrupan en dos enfoques principales (ver Figura \ref{fig:MI1}): la \textbf{selección de características}, que conserva un subconjunto de las variables originales mediante métodos de filtro o enfoques \textit{wrapper}; y la \textbf{extracción de características}, que transforma el espacio original para generar nuevas proyecciones, como ocurre en el Análisis de Componentes Principales (PCA) o el Análisis Discriminante \cite{jain2000statistical}.

\begin{figure}[htbp]
	\centering
	\includegraphics[width=0.7\textwidth]{D1.png}
	\caption{Taxonomía de las técnicas de reducción de dimensionalidad. Tomada de \cite{medina2025genetic}.}
	\label{fig:MI1}	
\end{figure}

La proyección de los datos hacia espacios de baja dimensión (dos o tres dimensiones) cumple un propósito exploratorio fundamental. En dominios complejos como la genómica o la radiómica, visualizar la geometría de los datos permite evaluar si los biomarcadores extraídos logran discriminar claramente entre condiciones sanas y patológicas \cite{bruno2020data, neha2026radiomics}. Comúnmente, esta exploración se realiza proyectando las dos primeras componentes de PCA \cite{kabir2023performance}.

\section{Planteamiento del problema}

A pesar de su popularidad, PCA no es la herramienta más adecuada para evaluar la separabilidad de clases. Al ser un método no supervisado cuyo criterio de optimización es la varianza total, no garantiza que las proyecciones obtenidas coincidan con las fronteras que mejor separan a los grupos \cite{rabbi2025performance}. El Análisis Discriminante Lineal (LDA) resuelve esta limitación incorporando la información de las etiquetas para maximizar la dispersión inter-clase respecto a la intra-clase, convirtiéndolo en la herramienta natural para la visualización supervisada \cite{chhikara2022data}.

No obstante, la formulación clásica de LDA impone supuestos que restringen severamente su aplicabilidad: asume fronteras de decisión estrictamente lineales, normalidad multivariada y matrices de dispersión bien condicionadas, fallando rotundamente en escenarios con tamaños de muestra pequeños \cite{feng2021comparison}. 

Para superar estas barreras operativas, se han desarrollado variantes evolutivas como el Genetic Projective Discriminant Analysis (GPDA), el cual optimiza la proyección discriminante lineal esquivando los problemas de matrices singulares \cite{medina2025genetic}. Sin embargo, dicha propuesta conserva la restricción de su método base: su formulación esta diseñada para tratar con datos con estructuras linealmente separables.

Una alternativa intuitiva para abordar la no linealidad es expandir explícitamente el espacio de características (por ejemplo, mediante términos polinomiales), convirtiendo relaciones complejas en fronteras lineales tratables por LDA. No obstante, este enfoque introduce una explosión combinatoria del número de variables a medida que aumentan el grado del polinomio, agravando la maldición de la dimensionalidad \cite{bishop2006pattern}.

En la literatura actual, la alternativa dominante para tratar con estructuras no linealmente separables en los datos sin expandir explícitamente el espacio es el uso de métodos kernel, como las Maquinas de Soporte Vectorial (SVM) o el Análisis discriminante de Kernel Fisher (KFDA). Estos explotan el ``truco del kernel'' para operar implícitamente en espacios de dimensión alta o infinita. En particular, el KFDA realiza un mapeo de los datos hacia un espacio de mayor dimensionalidad y, posteriormente, los proyecta en una subdimensión reducida de tamaño  ($C - 1$) número de clases.Este proceso de reducción se lleva a cabo mediante combinaciones matemáticas de los datos transformados, lo que posibilita la visualización en baja dimensionalidad, pero dificulta significativamente la interpretación directa de la representación en términos de las variables originales. 

Existe, por tanto, una brecha metodológica concreta: no se cuenta con un mecanismo capaz de producir una proyección de baja dimensión para estructuras no linealmente separables que sea simultáneamente visualizable y trazable hacia las variables originales. 

\section{Propuesta}

El objetivo es diseñar e implementar un método híbrido que aplique una expansión polinomial explícita para modelar fronteras de decisión no lineales. Se utilizará dentro de un algoritmo genético jerárquico de dos etapas como esquema de envoltura para seleccionar los términos óptimos, empleando el Análisis Discriminante Lineal (LDA) como función de evaluación y motor para proyectar los datos en un subespacio de máxima separabilidad visualizable en máximo tres dimensiones.

\section{Justificación}
\label{sec:Justificacion}

Esta investigación se justifica por la necesidad de encontrar un equilibrio óptimo entre la transformación para la separabilidad no lineal y la parsimonia dimensional. El trabajo propuesto desarrolla un algoritmo de tipo \textit{wrapper} capaz de solventar la no linealidad mediante expansión polinomial, mitigando simultáneamente la maldición de la dimensionalidad a través de una selección evolutiva de variables.

A diferencia de métodos de ``caja negra'' del estado del arte, como el Análisis Discriminante de Fisher con Kernel (KFDA) o las Máquinas de Soporte Vectorial (SVM), este marco de trabajo ofrece un beneficio dual:

\begin{itemize}
	\item \textbf{Visualización geométrica:} Permite proyectar estructuras no linealmente separables en espacios de baja dimensión (hasta tres dimensiones), facilitando la inspección visual de la separación entre clases.
	
	\item \textbf{Trazabilidad e interpretabilidad:} El proceso permite rastrear con precisión qué variables originales y qué términos de interacción fueron seleccionados, ofreciendo una explicación técnica de la proyección final. Esto resulta de suma importancia en aplicaciones donde entender el ``porqué'' de la decisión del modelo es tan importante como la precisión misma.
\end{itemize}

\section{Pregunta de investigación}
\label{sec:Pregunta-Investigación}

A partir de esta problemática, la presente investigación plantea la siguiente pregunta:

\begin{itemize}
	\item \textbf{¿Es posible integrar la expansión polinomial explícita con un esquema de selección evolutiva en dos fases para proyectar datos con estructuras no linealmente separables en espacios de baja dimensión (hasta tres dimensiones), que sean visualizables, interpretables y trazables a partir de las variables y términos seleccionados, manteniendo al mismo tiempo un desempeño predictivo y un costo computacional competitivos frente a métodos no lineales de referencia?}
\end{itemize}

\section{Hipótesis}
\label{sec:Hipótesis}

La integración de un mecanismo de selección evolutiva en dos fases sobre una expansión polinomial explícita permite generar representaciones discriminantes visualizables en baja dimensión (hasta tres dimensiones) para conjuntos de datos con estructuras no linealmente separables. El enfoque permite reducir la dimensionalidad del espacio expandido e identificar los términos polinomiales seleccionados que contribuyen a la separación discriminante, proporcionando trazabilidad sobre el proceso de construcción de la representación y favoreciendo su interpretabilidad. Asimismo, se plantea que esta representación puede obtenerse manteniendo tiempos de ejecución viables y una exactitud de clasificación estadísticamente competitiva frente a métodos no lineales de referencia de tipo caja negra, sin requerir una superioridad sistemática en rendimiento predictivo.

\section{Objetivos}
\label{sec:Objetivo-General}

\subsection{Objetivo general}

Diseñar e implementar un algoritmo con un esquema híbrido de selección y extracción de características que integre expansión polinomial con análisis discriminante lineal mediante un proceso evolutivo en dos etapas, con el propósito de lograr la visualización interpretable y trazable de datos con estructuras no linealmente separables en espacios de baja dimensión, manteniendo un equilibrio entre reducción de características, tiempo de ejecución y rendimiento predictivo.

\subsection{Objetivos específicos}

\begin{itemize}
	
	\item \textbf{Caracterizar} las limitaciones operativas y visuales del Análisis Discriminante Lineal (LDA) clásico ante escenarios de no linealidad, alta dimensionalidad y muestras pequeñas.
	
	\item \textbf{Diseñar e implementar} un mecanismo de selección evolutiva que reduzca la redundancia en el espacio expandido mediante la selección de variables originales y términos polinomiales relevantes que sea capaz de generar una proyección en baja dimensión trazable.
	
	\item \textbf{Evaluar} el desempeño computacional y la reducción dimensional producida por el algoritmo de selección evolutiva, mediante el registro del tiempo de ejecución y la cuantificación de la cantidad de variables retenidas en cada una de sus fases.	
	
	\item \textbf{Comparar de forma cualitativa} las proyecciones visuales de baja dimensión (máximo tres dimensiones) generadas por el método propuesto frente a las representaciones producidas por LDA y KFDA, evaluando su utilidad para la exploración e interpretación de estructuras discriminantes no lineales.		
	
	\item \textbf{Evaluar y contrastar} estadísticamente la exactitud de clasificación (\textit{Accuracy}) de la propuesta frente a las técnicas de referencia, para validar su competitividad predictiva.
	
	%\item \textbf{Analizar} estadísticamente las diferencias de desempeño entre los %métodos evaluados.
	
\end{itemize}

\section{Contribuciones}
\label{sec:Contribuciones}

Las principales contribuciones de este trabajo son las siguientes:

\begin{itemize}
	
	\item Proponer un esquema híbrido de extracción de características que combina expansión polinomial, selección evolutiva y análisis discriminante para abordar problemas de
	clasificación y proyección con datos no linealmente separables.
	
	\item Incorporar un mecanismo de selección en dos fases orientado a reducir la dimensionalidad del espacio expandido, primero mediante la selección de variables originales y posteriormente mediante la selección de términos polinomiales relevantes.
	
	\item Evaluar la capacidad del método propuesto para generar proyecciones de baja dimensión, hasta un máximo de tres dimensiones que faciliten la visualización de la separación entre clases.
	
	\item Demostrar empíricamente que el algoritmo propuesto alcanza una exactitud de clasificación estadísticamente competitiva frente a métodos que manejan la no linealidad (modelos de caja negra), aportando como ventaja fundamental una proyección interpretable y trazable sin incurrir en una pérdida significativa de rendimiento predictivo.
	
	\item Cuantificar la eficiencia del mecanismo evolutivo de selección en sus dos fases mediante métricas tangibles de costo computacional (tiempo de ejecución) y optimización de memoria (tasa de reducción de variables del espacio polinomial expandido), validando su viabilidad operativa.
	
\end{itemize}

\section{Alcances y limitaciones}
\label{sec:Alcances-Limitaciones}


\begin{itemize}
	
	\item El enfoque de la investigación se restringe a problemas de aprendizaje supervisado.
	
	\item El alcance experimental se restringe de forma estricta a conjuntos de datos tabulares donde las características exclusivamente numéricas. Por tanto, no se consideran en esta etapa datos no estructurados, variables categóricas sin transformación previa, imágenes, texto, señales temporales ni otros tipos de representación.
	
	\item Las visualizaciones generadas por el método propuesto se acotan a espacios de baja dimensión, hasta un máximo de tres dimensiones. Esta restricción se relaciona con la naturaleza del análisis discriminante lineal, el cual permite obtener hasta $C - 1$ dimensiones discriminantes, donde $C$ representa el número de clases del problema.
	
	\item La propuesta no busca superar en rendimiento a todos los métodos de referencia, sino obtener un desempeño comparable en determinados escenarios, aportando como ventaja la trazabilidad del proceso de selección y la posibilidad de visualizar las proyecciones generadas en baja dimensión.
	
	\item La validación empírica del algoritmo se restringe a un conjunto controlado de diez conjuntos de datos, con un máximo de 32 variables cuantitativas originales.
	
\end{itemize}

\section{Organización de la tesis}
\label{sec:Organización-Tesis}

\begin{itemize}
	
	\item El \textbf{Capítulo 2} presenta el marco teórico necesario para comprender los fundamentos de la investigación. En este capítulo se abordan las técnicas de reducción de dimensionalidad, los fundamentos algebraicos y operacionales del Análisis Discriminante Lineal y sus variantes, el problema de pequeño tamaño de muestra, la expansión polinomial, los algoritmos genéticos, los métodos kernel, KFDA, las máquinas de soporte vectorial y las métricas de clasificación utilizadas en el estudio.
	
	\item El \textbf{Capítulo 3} desarrolla el estado del arte. En este capítulo se revisan trabajos relacionados con variantes y mejoras del Análisis Discriminante Lineal, métodos kernel como KPCA, KFDA y SVM, así como enfoques de selección de características basados en algoritmos genéticos. Asimismo, se analizan trabajos relacionados con GPDA y se identifican las limitaciones que motivan el desarrollo de la propuesta.
	
	\item El \textbf{Capítulo 4} describe la metodología de investigación y el algoritmo propuesto, denominado Genetic Non-Linear Polynomial Discriminant Analysis (GNLPDA). Se presentan la arquitectura general del método, sus fases de selección, la función fitness, el proceso de calibración y el diseño experimental empleado para su evaluación.
	
	\item El \textbf{Capítulo 5} presenta los resultados obtenidos al aplicar el método propuesto sobre los conjuntos de datos seleccionados. Se incluyen resultados de clasificación, análisis de proyecciones de baja dimensión y comparaciones frente a los métodos de referencia.
	
	\item El \textbf{Capítulo 6} discute los hallazgos obtenidos, considerando la interpretación técnica de los resultados, las fortalezas del método, sus limitaciones y las condiciones bajo las cuales la propuesta resulta más adecuada.
	
	\item El \textbf{Capítulo 7} presenta las conclusiones generales de la investigación, el grado de cumplimiento de los objetivos planteados, las principales contribuciones del trabajo y las líneas de trabajo futuro.
	
\end{itemize}




